\section{Extended robustness and route results}
\label{app:extended}

\subsection{Post-confirmatory falsification}

The preferred explanation survived in qualified form. Four bootstrap seeds preserve the selected-versus-random safety gate. No leave-one-out deletion changes support; after removing the most favorable 20\% of examples, the minimum frozen gate value remains 0.693. Mean, trimmed-mean, and five-group median-of-means aggregations agree. Twenty thousand sign-flip max-$T$ permutations across the frozen 68-candidate Day 8 family retain all 16 selected candidates at familywise 0.05.

Several strict hypotheses fail. Removing the selected MLPs leaves nearly the full effect, rejecting minimal mixed-K16 wording. Type-compatible site shuffling preserves strong induction under some mappings, rejecting rigid source-to-original-site necessity. Zero ablation recovers much less than natural-state rescue, rejecting it as an adequate substitute. Harmfulness induction is null at one absolute threshold because both compared conditions saturate, rejecting universal hard-threshold wording.

\subsection{Route-population confirmation}

\begin{table}[h]
\centering
\caption{Prospective Day 21 population effects. Values are estimate [95\% paired-bootstrap interval].}
\label{tab:route-gate}
\begin{tabular}{llll}
\toprule
Concept & Direction & Selected delta & Selected $-$ matched null \\
\midrule
Deception & Rescue & 0.133 [0.114, 0.153] & 0.134 [0.115, 0.153] \\
Deception & Induction & 0.387 [0.362, 0.410] & 0.388 [0.363, 0.412] \\
Harmfulness & Rescue & 0.164 [0.151, 0.177] & 0.190 [0.175, 0.204] \\
Harmfulness & Induction & 0.093 [0.082, 0.104] & 0.102 [0.090, 0.114] \\
\bottomrule
\end{tabular}
\end{table}

The bootstrap samples examples with the four mappings fixed. It therefore quantifies example uncertainty, not uncertainty over possible mapping generators or receiver populations. Mapping-level heterogeneity is part of the result rather than a nuisance averaged away.

\subsection{Behavioral transport disposition}

\begin{table}[h]
\centering
\caption{Frozen Day 23 population endpoints. Probe and directional values include 95\% intervals; KL and NLL are point estimates in nats/token.}
\label{tab:behavior-gate}
\begin{tabular}{llllll}
\toprule
Concept & Direction & Probe transport & Output coefficient & KL & NLL shift \\
\midrule
Deception & Induction & 0.386 [0.366, 0.405] & 0.022 [0.019, 0.026] & 0.0025 & $-0.0042$ \\
Deception & Rescue & 0.138 [0.122, 0.154] & 0.028 [0.024, 0.032] & 0.0029 & 0.0017 \\
Harmfulness & Induction & 0.090 [0.075, 0.103] & 0.087 [0.080, 0.094] & 0.0098 & 0.0188 \\
Harmfulness & Rescue & 0.161 [0.146, 0.175] & 0.098 [0.090, 0.105] & 0.0130 & $-0.0195$ \\
\bottomrule
\end{tabular}
\end{table}

The mechanical labels are behavior-preserving portable evasion for both deception directions and harmfulness induction, and mixed for harmfulness rescue. These labels mean that the specified KL/NLL equivalence checks pass and the population coefficient is below or boundary-mixed at the frozen 0.10 scale. They do not imply exactly unchanged logits, strings, semantics, truthfulness, or safety.

\subsection{Package integrity}

The Day 13 canonical archive contains exactly 1,944 intervention rows. The Day 23 teacher-forced behavioral archive contains 3,880 rows; the Day 24 coupled-generation archive contains 80 rows. All have recorded SHA-256 digests, exact condition-grid audits, finite-value checks, hook cleanup checks, and Git ancestry linking implementation, authorization, and raw results. The final local test suite passes 62 tests. No independent evaluator or external replication is included.
